% !TEX root = ../main.tex
\section{相关工作}

\subsection{复杂网络与网络动力学}

复杂网络研究关注连接结构如何影响系统的集体行为。
随机网络、小世界网络和无标度网络分别提供了均匀随机连接、局部聚集与短路径并存、度分布高度异质等典型结构\cite{erdos1959,watts1998,barabasi1999}。
在这些网络上，
节点状态并非独立变化，
而是由自身规律和局部耦合共同驱动；
相同的节点动力学置于不同拓扑上，
可能呈现不同的传播阈值、同步模式和稳定性\cite{strogatz2001,albert2002}。

网络 SIR 是研究“结构如何塑造演化”的经典模型。
相较于均匀混合假设，网络模型把传播限定在连接边上，因此能够表达聚类、枢纽节点和局部暴发等效应\cite{pastorsatorras2015}。
本文使用 ER、WS 和 BA 网络并非为了把三者当作现实系统的替代，
而是希望在连接均匀性、局部聚集性和度异质性明显不同的条件下检验学习方法。

\subsection{数据驱动动力学学习}

数据驱动的动力学发现试图从观测轨迹中恢复状态演化规律。
稀疏辨识等方法通过预设函数库寻找简洁方程，为“从数据到规律”提供了重要思路\cite{brunton2016}；
神经网络方法则以更强的函数逼近能力处理难以事先列举的非线性关系。
对网络化系统而言，
核心并不是简单地把所有节点状态拼接为一个高维输入，
而是利用交互的局部性和共享性：每个节点遵循相似的更新形式，
但接收到的邻居作用由网络决定。

图神经网络恰好为这类关系提供了自然表示。
图卷积和 GraphSAGE 等方法通过聚合邻居信息完成节点级预测\cite{kipf2017gcn,hamilton2017}；
GIN 的分析进一步说明，
聚合函数能否保留邻居多重集合信息会影响模型的区分能力\cite{xu2019gin}。
本文使用的 GraphConv 基线对自身项与邻居项分别变换后相加\cite{morris2019,fey2019}。
在此基础上，
我们进一步保持两类表示的分离，
使模型可以独立调整自身演化和交互作用，
而不是把双通道优势归结为基线在结构上的“必然不可区分”。

\subsection{潜在关系学习与网络重构}

当拓扑不可见时，动力学学习还要回答“谁与谁发生作用”。
Neural Relational Inference 从多主体轨迹推断离散关系类型\cite{kipf2018nri}；
Gumbel Graph Network（GGN）则将网络生成器与动力学学习器联合训练，
以离散关系的连续松弛实现端到端优化\cite{zhang2019ggn}。
Gumbel-Softmax 和 Concrete 分布为这类离散采样提供了可微近似\cite{jang2017,maddison2017}。

联合学习的关键风险在于结构与动力学互相补偿：错误连接不一定立即造成较大预测误差，灵活的动力学模型可能通过改变交互强度吸收结构偏差。
本文沿用“生成网络并用预测误差反向约束结构”的基本思想，
但将第一阶段得到的双通道表示引入未知拓扑任务，
同时在硬图上学习规律，
并让结构更新使用软图梯度。
研究重点因此不是单纯提高一步预测精度，
而是在预测、规律诊断和边级恢复之间建立相互校验。

\subsection{本文研究定位}

已有研究分别提供了数据驱动方程发现、图表示学习和潜在关系推断的基础。
本文将这些线索组织为一个递进问题：先在网络已知时弄清怎样学习可解释、可推演的局部演化规律，
再在网络未知时处理规律与结构的共同辨识。
SIR 是本文用于验证这一路径的具体动力学，
而非方法成立的唯一对象。
方法所依赖的更一般假设是：组分状态可观测，
系统在给定时间尺度上近似满足局部马尔可夫更新，
交互能够由邻居状态的某种聚合表示。
